\documentclass[10pt]{article}
\usepackage[margin=1in]{geometry}
\usepackage{amsmath}
\usepackage{booktabs}
\usepackage{longtable}
\usepackage{array}
\usepackage{placeins}
\usepackage{hyperref}

\title{Auxiliary Check: EE-D at $N{=}100$ vs.\ the Paper's $N{=}10$ Numbers\\for the RandMMR-vs-PL Comparison}
\author{}
\date{}

\begin{document}
\maketitle

\noindent\textbf{This is a working diagnostic document, not part of the paper.} It
checks how much Table 2 / Appendix B's per-cell EE-D numbers would move if EE-D were
estimated from $N{=}100$ reranking samples instead of $N{=}10$, for the same
comparison the paper already reports: the (buggy) hybrid reranker
\texttt{pl\_randmmr} ($\alpha{=}4$, $\lambda\sim U(0.8,1.0)$) against its matched
plain-PL setting, per (task, generator, ranker) cell. EE-D only depends on which
documents land in which ranked positions across the sampled lists, not on any
generated text, so this recomputation needed zero LLM calls --- retrieval and
reranking only. Expected Utility (EU$_{\text{norm}}$) and its significance ($p$-value)
are left exactly as already reported from the real $N{=}10$ generation runs; only
EE-D is recomputed at higher precision.

\section{Why this check}

Every EE-D number anywhere in this project is a Monte Carlo estimate over a finite
set of $N$ sampled rankings. A separate check (not shown here) found that low-entropy
policies (e.g.\ high-$\alpha$ PL) barely move between $N{=}10$ and $N{=}100$, while
high-entropy policies (e.g.\ low-$\alpha$ PL, or the stochastic-tail hybrid reranker)
can shift substantially --- $N{=}10$ systematically \emph{overstates} disparity for
these because a small sample under-explores how spread out the policy really is. Since
the paper's Table 2 comparison uses $N{=}10$ throughout, it is worth checking directly
whether the reported RandMMR-vs-PL gap survives a higher-precision re-estimate, rather
than assuming it does.

\section{Result}

\noindent\textbf{Which PL setting.} The paper matches each cell's hybrid reranker
against whichever PL$(\alpha)$ has the closest own EE-D (biased toward equal-or-higher
PL disparity when no close match exists; see the main paper's methodology notes for
the exact rule). That per-cell match happens to land on $\alpha{=}2$ for \emph{every}
one of the 28 cells in this particular comparison --- not assumed, checked directly
against the \texttt{chosen\_pl\_alpha} column of
\texttt{rq3\_randmmr\_vs\_pl\_n10\_all\_cells.csv}. So every ``PL EE-D'' value below is
PL$(\alpha{=}2)$, not a mix of different $\alpha$'s across rows.

\textbf{It survives, and if anything strengthens slightly.} RandMMR has the lower
(fairer) EE-D in 26 of 28 cells at both $N{=}10$ and $N{=}100$. The mean gap
(RandMMR $-$ PL) goes from $-0.0089$ at $N{=}10$ to $-0.0113$ at $N{=}100$ --- wider in
RandMMR's favor, not narrower. Four individual cells flip direction (marked
$^\dagger$ below), but they cancel out: two flip from PL-favoring to RandMMR-favoring
(both essentially exact ties at $N{=}10$, margins under $0.002$), and two flip the
other way (LaMP-7, both BM25 cells, margins under $0.003$ either direction) --- none
of the four flips is a real reversal, all are noise-level margins at $N{=}10$ resolving
in one direction or the other once precision improves. Utility is unaffected throughout:
the same two cells remain significant (LaMP-2 Small/BM25, $p{=}0.017$; LaMP-4
Small/BM25, $p{=}0.0019$), both already favoring RandMMR on utility.

\FloatBarrier
\begin{longtable}{c l l cc cc rr}
\toprule
LaMP & Gen. & Ranker & \multicolumn{2}{c}{PL($\alpha{=}2$) EE-D} & \multicolumn{2}{c}{RandMMR EE-D} & \multicolumn{2}{c}{$\Delta$EE-D (RandMMR$-$PL)} \\
 & & & $N{=}10$ & $N{=}100$ & $N{=}10$ & $N{=}100$ & $N{=}10$ & $N{=}100$ \\
\midrule
\endfirsthead
\toprule
LaMP & Gen. & Ranker & \multicolumn{2}{c}{PL($\alpha{=}2$) EE-D} & \multicolumn{2}{c}{RandMMR EE-D} & \multicolumn{2}{c}{$\Delta$EE-D (RandMMR$-$PL)} \\
 & & & $N{=}10$ & $N{=}100$ & $N{=}10$ & $N{=}100$ & $N{=}10$ & $N{=}100$ \\
\midrule
\endhead
\bottomrule
\endlastfoot
1 & Base & BM25 & 0.193 & 0.115 & 0.186 & 0.099 & -0.0070 & -0.0157 \\
1 & Base & Contr. & 0.182 & 0.100 & 0.173 & 0.088 & -0.0087 & -0.0122 \\
1 & Small & BM25 & 0.178 & 0.098 & 0.173 & 0.084 & -0.0055 & -0.0138 \\
1 & Small & Contr. & 0.172 & 0.088 & 0.163 & 0.077 & -0.0091 & -0.0112 \\
2 & Base & BM25 & 0.378 & 0.316 & 0.356 & 0.302 & -0.0222 & -0.0140 \\
2 & Base & Contr. & 0.378 & 0.315 & 0.358 & 0.302 & -0.0204 & -0.0129 \\
2 & Small & BM25 & 0.338 & 0.292 & 0.338 & 0.277 & +0.0000 & -0.0151$^\dagger$ \\
2 & Small & Contr. & 0.357 & 0.291 & 0.338 & 0.277 & -0.0184 & -0.0137 \\
3 & Base & BM25 & 0.151 & 0.065 & 0.146 & 0.057 & -0.0052 & -0.0086 \\
3 & Base & Contr. & 0.150 & 0.064 & 0.144 & 0.055 & -0.0055 & -0.0085 \\
3 & Small & BM25 & 0.143 & 0.062 & 0.144 & 0.055 & +0.0012 & -0.0079$^\dagger$ \\
3 & Small & Contr. & 0.148 & 0.062 & 0.145 & 0.054 & -0.0039 & -0.0080 \\
4 & Base & BM25 & 0.254 & 0.183 & 0.251 & 0.174 & -0.0032 & -0.0085 \\
4 & Base & Contr. & 0.245 & 0.173 & 0.238 & 0.163 & -0.0068 & -0.0101 \\
4 & Small & BM25 & 0.259 & 0.177 & 0.246 & 0.169 & -0.0127 & -0.0087 \\
4 & Small & Contr. & 0.239 & 0.167 & 0.233 & 0.157 & -0.0051 & -0.0102 \\
5 & Base & BM25 & 0.181 & 0.103 & 0.173 & 0.089 & -0.0077 & -0.0138 \\
5 & Base & Contr. & 0.184 & 0.106 & 0.178 & 0.092 & -0.0061 & -0.0141 \\
5 & Small & BM25 & 0.180 & 0.103 & 0.174 & 0.089 & -0.0064 & -0.0140 \\
5 & Small & Contr. & 0.184 & 0.107 & 0.179 & 0.093 & -0.0046 & -0.0142 \\
6 & Base & BM25 & 0.239 & 0.164 & 0.231 & 0.149 & -0.0085 & -0.0152 \\
6 & Base & Contr. & 0.236 & 0.158 & 0.226 & 0.144 & -0.0101 & -0.0138 \\
6 & Small & BM25 & 0.237 & 0.163 & 0.230 & 0.148 & -0.0072 & -0.0154 \\
6 & Small & Contr. & 0.235 & 0.157 & 0.226 & 0.144 & -0.0090 & -0.0136 \\
7 & Base & BM25 & 0.505 & 0.465 & 0.499 & 0.468 & -0.0058 & +0.0032$^\dagger$ \\
7 & Base & Contr. & 0.488 & 0.437 & 0.463 & 0.423 & -0.0242 & -0.0138 \\
7 & Small & BM25 & 0.535 & 0.501 & 0.530 & 0.502 & -0.0047 & +0.0011$^\dagger$ \\
7 & Small & Contr. & 0.523 & 0.479 & 0.500 & 0.466 & -0.0231 & -0.0132 \\

\end{longtable}
\noindent $^\dagger$ cell where the RandMMR-vs-PL EE-D comparison flips direction between $N{=}10$ and $N{=}100$ (all four are noise-level margins at $N{=}10$, under $0.003$).

\section{Caveat}

Decoupling EE-D's sample size from EU's is methodologically valid --- EE-D/EE-R/EE-L
are computed purely from reranking output via \texttt{compute\_ee}, with no
dependence on which lists were actually sent to the LLM, so EE-D at $N{=}100$ and EU
at $N{=}10$ both remain correct Monte Carlo estimates of the same policy's true EE-D
and true EU respectively. The one thing given up is that EE-D and EU no longer
describe the identical realized set of served lists (only the identical
\emph{policy}) --- irrelevant for the policy-level claims this project makes, but
worth naming explicitly.

\end{document}
